\chapter{Orchestration and Reliability}
\label{chapter:Orchestration}
% EVOLVE-BLOCK-START

A single agent executing a single task is the simplest case of agentic work.\index{Orchestration}
Real research tasks are more demanding: they require long, multi-step computations; they benefit from having one agent check another's work; they must run reliably for hours without supervision; and they must report failures honestly when they do fail.
\emph{Orchestration} is the practice of composing multiple agents and operations into a trustworthy workflow.

Challenge~4 builds an orchestrated pipeline that previews the evolutionary framework.
The loop implemented here — \textbf{propose → evaluate → select → repeat} — is the same loop that becomes an evolutionary algorithm when a mutation operator is added at the pivot.
Section~\ref{section:MultiStepWorkflows} composes single operations into pipelines; Section~\ref{section:MultiAgentArchitectures} adds a second agent to check the first; Section~\ref{section:MemoryState} handles what must be remembered across steps; Section~\ref{section:Reliability} makes the pipeline survive long unattended runs; and Section~\ref{section:GenEvalSel} distills all of it into the generate--evaluate--select loop that the next chapter turns into evolution.

\section{Multi-Step Workflows}
\label{section:MultiStepWorkflows}

A \emph{workflow}\index{Workflow} is a directed graph of operations: some operations depend on the outputs of others; some can run in parallel.
The orchestrator's job is to schedule these operations correctly and collect their results.

\subsection{Pipeline Composition}

In the simplest case, a workflow is a \emph{pipeline}\index{Pipeline}: a linear sequence of operations where the output of each step is the input to the next.

\begin{example}
A three-step research pipeline:
\begin{lstlisting}[language=python]
def run_pipeline(dataset_path, k=5):
    # Step 1: retrieve domain knowledge
    context = retriever.query(dataset_path, k=k)

    # Step 2: agent proposes and scores candidates
    candidates = proposer.run(dataset_path, context=context)

    # Step 3: select the best candidate
    best = min(candidates, key=lambda c: c.score)
    return best
\end{lstlisting}
Each step can be a model call, a tool invocation, or a subprocess; the pipeline connects them.
\end{example}

\subsection{Parallel Execution}

Some operations are independent: multiple agents can propose candidates simultaneously, or a critic can score candidates in parallel with new proposals being generated.
\emph{Parallel execution}\index{Parallel execution} reduces wall-clock time by running independent operations concurrently.

\begin{example}
Spawning multiple proposer agents in parallel (conceptual):
\begin{lstlisting}[language=python]
import concurrent.futures

with concurrent.futures.ThreadPoolExecutor(max_workers=4) as pool:
    futures = [pool.submit(proposer.run, dataset_path, seed=i)
               for i in range(4)]
    all_candidates = [f.result() for f in futures]
\end{lstlisting}
The four proposers run concurrently; the orchestrator collects and merges their outputs.
\end{example}

The principal hazard of parallel execution is \emph{non-determinism}\index{Non-determinism}: the order in which parallel operations complete is not predictable across runs.
If the final selection depends on ordering (e.g., in case of a tie), the result may differ between runs.
Always set seeds and record exactly which candidates were selected.

\section{Multi-Agent Architectures}
\label{section:MultiAgentArchitectures}

\subsection{The Proposer--Critic Pattern}

The most common multi-agent pattern in the course is the \emph{proposer--critic} pattern (Figure~\ref{figure:ProposerCritic}):\index{Proposer--critic pattern}
\begin{itemize}
\item A \emph{proposer} agent generates candidate artifacts: model expressions, hypotheses, experimental designs.
Its goal is breadth — generating diverse, plausible candidates quickly.
\item A \emph{critic} (or \emph{verifier}) agent examines each candidate for correctness, plausibility, or constraint satisfaction before the candidate is passed to the evaluator.
Its goal is precision — filtering out candidates that are obviously wrong.
\end{itemize}

\begin{figure}[htb!]
\begin{center}
\begin{tikzpicture}[
    node distance=1.5cm and 1.8cm,
    box/.style={draw, rounded corners=4pt, minimum width=2.4cm,
                minimum height=0.8cm, align=center, font=\small},
    arr/.style={->, >=stealth', thick}
]
\node[box, fill=blue!15]   (prop) {Proposer};
\node[box, fill=orange!20, right=of prop] (crit) {Critic};
\node[box, fill=green!20,  right=of crit] (eval) {Evaluator};
\node[box, fill=gray!15,   below=of eval] (sel)  {Selector};

\draw[arr] (prop) -- node[above,font=\scriptsize]{candidates} (crit);
\draw[arr] (crit) -- node[above,font=\scriptsize]{survivors}  (eval);
\draw[arr] (eval) -- node[right,font=\scriptsize]{scores}     (sel);
\draw[arr] (sel.west) -- ++(-3.8,0) |-
    node[pos=0.25,below,font=\scriptsize]{iterate / report} (prop.west);
\end{tikzpicture}
\caption{The proposer--critic--evaluator architecture.
The proposer generates candidates, the critic filters out obvious errors, the evaluator scores the survivors, and the selector feeds the best back to the proposer for the next round.}
\label{figure:ProposerCritic}
\end{center}
\end{figure}

The critic's role is to enforce constraints that the evaluation harness does not check.
In the symbolic regression context, a critic might check:
\begin{itemize}
\item dimensional consistency of a proposed expression,
\item whether the expression is syntactically valid Python,
\item whether the expression is numerically stable on the training domain, and
\item whether the expression has already been tried in a previous round (deduplication).
\end{itemize}
Filtering before evaluation saves evaluation cost — especially important when evaluation requires model calls or expensive computation.
This proposer--critic loop is the pattern Anthropic's engineering writing (2024) names the \emph{evaluator--optimizer} workflow, and it presupposes what Chapter~\ref{chapter:Evaluation} supplies: explicit criteria against which iteration demonstrably improves the artifact, which is why the critic is a filter and not a rubber stamp.

\subsection{Verifiers}

A \emph{verifier}\index{Verifier} is an agent whose explicit task is to try to \emph{disprove} a claim or detect an error.
Where a critic checks soft constraints, a verifier checks hard correctness properties.

\begin{example}
A verifier for a data analysis step:
\begin{lstlisting}[language=markdown]
You are a verification agent. The primary agent has made the following
claim about the dataset:

  "The features are all non-negative, so log-transforms are safe."

Your job is to check this claim against the actual data. Run the
following check and report whether the claim is true or false:

  import numpy as np
  X, y = load_dataset('data/tier1_growth.npz')
  print(np.min(X), np.any(X <= 0))
\end{lstlisting}
The verifier can catch mistakes the proposer made silently — here, the possibility that some feature values are zero or negative.
\end{example}

\subsection{Dynamic Sub-Agent Spawning}

The proposer--critic topology above is \emph{static}: the agents and their wiring are fixed before the run begins.
A more flexible pattern lets an \emph{orchestrator}\index{Orchestrator agent} agent spawn specialist sub-agents on demand --- calling a tool that launches a fresh agent instance with its own context window, tool set, and task, then collecting the result.
This is how a single agent decomposes a large job at run time: a research orchestrator might spawn one sub-agent per candidate dataset, or one per section of a report, each working in an isolated context so that the orchestrator's own window is not flooded with intermediate detail (Figure~\ref{figure:OrchestratorWorker}).
The reliability discipline of this chapter applies unchanged --- each spawned agent can fail, so its result must be validated and its failures surfaced (Section~\ref{section:Reliability}).
Current agent SDKs (for example, Claude's agent SDK) expose this spawning as a built-in tool; it is the dynamic generalization of the fixed proposer--critic wiring.
When parallel agents must \emph{write} --- not just analyze --- they also need isolated working directories so their changes cannot collide; the git-worktree pattern that provides this is described in Appendix~\ref{chapter:RigReference}.

\begin{figure}[htb!]
\begin{center}
\begin{tikzpicture}[
    box/.style={draw, rounded corners=4pt, minimum width=2.3cm, minimum height=1.0cm, align=center, font=\small},
    sub/.style={draw, rounded corners=4pt, minimum width=2.7cm, minimum height=0.9cm, align=center, font=\scriptsize},
    arr/.style={->, >=stealth', thick}
]
\node[box, fill=orange!20] (orch) at (0,0) {Orchestrator};
\node[sub, fill=blue!12] (s1) at (4.7,1.7)  {sub-agent\\\emph{own context}};
\node[sub, fill=blue!12] (s2) at (4.7,0)    {sub-agent\\\emph{own context}};
\node[sub, fill=blue!12] (s3) at (4.7,-1.7) {sub-agent\\\emph{own context}};
\node[box, fill=green!20] (syn) at (9.4,0) {Synthesis};
\draw[arr] (orch) -- (s1);
\draw[arr] (orch) -- (s2) node[midway, above, font=\scriptsize]{task};
\draw[arr] (orch) -- (s3);
\draw[arr] (s1) -- (syn);
\draw[arr] (s2) -- (syn) node[midway, above, font=\scriptsize]{result};
\draw[arr] (s3) -- (syn);
\end{tikzpicture}
\caption{The orchestrator--worker pattern (Anthropic, 2025).
A lead agent decomposes a task and spawns worker sub-agents that run in parallel, each with its own context window, then synthesizes their results.
Because each worker's context is isolated, the orchestrator's window is not flooded with intermediate detail and independent sub-questions are explored at once --- at the cost of roughly an order of magnitude more tokens than a single agent, so the topology pays off for breadth-first tasks whose parts are genuinely independent, not for tightly coupled work.}
\label{figure:OrchestratorWorker}
\end{center}
\end{figure}

\section{Memory and State}
\label{section:MemoryState}

A single-session agent remembers its history within the context window.
An orchestrated workflow that runs many rounds, or that spawns parallel agents, needs explicit \emph{external memory}\index{External memory}: state stored outside any individual agent's context.

\subsection{What Needs to Be Remembered}

\begin{itemize}
\item \textbf{Tried expressions.} The set of model forms already evaluated, so that the proposer does not waste budget on duplicates.
\item \textbf{Best results.} The highest-scoring candidates, maintained across rounds, so that the selector can always return the best discovered so far.
\item \textbf{Failure modes.} Candidates that crashed the evaluator (zero-division, overflow), along with the error, so the proposer can avoid similar forms.
\item \textbf{Run configuration.} The random seed, model name, dataset version, and harness settings, so the run can be reproduced exactly.
\end{itemize}

\subsection{Memory in Practice}

External memory is most simply a file — a JSON or CSV record updated after each round.

\begin{lstlisting}[language=python]
import json, pathlib

class RunMemory:
    def __init__(self, path):
        self.path = pathlib.Path(path)
        self.state = json.loads(self.path.read_text()) \
            if self.path.exists() else {"tried": [], "best": None}

    def record(self, expression, score):
        self.state["tried"].append({"expression": expression,
                                    "score": score})
        if (self.state["best"] is None
                or score < self.state["best"]["score"]):
            self.state["best"] = {"expression": expression,
                                  "score": score}
        self.path.write_text(json.dumps(self.state, indent=2))

    def already_tried(self, expression):
        return any(t["expression"] == expression
                   for t in self.state["tried"])
\end{lstlisting}
The memory file is the primary artifact of an orchestrated run: it records what was tried and what the best result was, enabling the run to be resumed if it is interrupted and reported honestly if it fails.

The single JSON file is the reproducible \emph{floor}, not the ceiling.
Production agents increasingly back external memory with a vector store, so that ``already tried'' can match \emph{semantically equivalent} candidates rather than only string-identical ones, and with hierarchical or episodic memory for runs that span many sessions (as of 2026).
The file-based version here is the minimal form that keeps a run resumable and honestly reportable; the richer stores are the same idea with a better index.

\section{Reliability Engineering}
\label{section:Reliability}

An orchestrated pipeline that runs unattended for hours must be designed to fail gracefully.
Every model call can return an error; every tool can produce an unexpected result; any network connection can drop.
\emph{Reliability engineering}\index{Reliability} is the practice of anticipating these failures and building the system to handle them without losing state or producing incorrect results.

\subsection{Retries with Backoff}

A \emph{retry}\index{Retry} re-attempts a failed operation after a waiting period.
\emph{Exponential backoff}\index{Exponential backoff} increases the waiting period after each successive failure, reducing the load on the server that produced the error.

\begin{lstlisting}[language=python]
import time, random
from anthropic import RateLimitError, APIError   # provider SDK exception types

def with_retry(fn, max_attempts=5, base_delay=1.0):
    for attempt in range(max_attempts):
        try:
            return fn()
        except (RateLimitError, APIError) as e:
            if attempt == max_attempts - 1:
                raise
            delay = base_delay * (2 ** attempt) + random.uniform(0, 1)
            time.sleep(delay)
    raise ValueError("with_retry needs max_attempts >= 1")
\end{lstlisting}

\subsection{Logging}

\emph{Logging}\index{Logging} records what the pipeline did, when, and with what results.
A run without logs is a run that cannot be understood after the fact.

Every orchestrated run should log:
\begin{itemize}
\item the start time and run configuration,
\item every model call (model name, token counts, latency),
\item every tool call and its result,
\item every error and its resolution, and
\item the final result and elapsed time.
\end{itemize}

\subsection{Honest Failure Reporting}

An orchestrated pipeline fails more often than a single-step script.
The response to failure is to \emph{surface it}, not to hide it.
A pipeline that silently returns a fallback when it fails is worse than one that raises an exception: the former produces incorrect results that look correct; the latter produces no result but an informative error.

\begin{principle}
\label{principle:HonestFailure}
Report failures explicitly.
If the pipeline does not find a candidate better than the linear baseline, say so.
If three of ten proposer calls failed due to rate limits, report the fraction.
If the critic rejected every proposal in a round, report why.
A failure report is not a sign of a bad system; a failure report that is absent is a sign of an untrustworthy one.
\end{principle}

\section{The Generate→Evaluate→Select Loop}
\label{section:GenEvalSel}

The orchestration assembled in Challenge~4 follows a pattern that recurs throughout the rest of the course.

\begin{definition}
\label{definition:GenEvalSel}
The \emph{generate–evaluate–select} (GES) loop is the procedure:
\begin{enumerate}
\item \textbf{Generate.} Produce a set of candidate artifacts $C = \{a_1, \ldots, a_m\}$.
\item \textbf{Evaluate.} Score each candidate: $s_i = \Score(a_i, \mathcal{D}_{\mathrm{test}})$.
\item \textbf{Select.} Retain the top-$k$ candidates: $C^* = \mathrm{top\text{-}}k(C, \{s_i\})$.
\item Repeat with $C^*$ informing the next round of generation.
\end{enumerate}
\end{definition}

Concretely: in round~1 the proposer generates five candidate expressions; the harness scores them $0.12$, $0.08$, $0.23$, $0.07$, $0.15$; with $k=2$ (and lower scores better), the selector retains the candidates scoring $0.07$ and $0.08$ as the elites that seed round~2.
Every run of the loop is this small story repeated.

This loop is an evolutionary algorithm in all but name.
The one ingredient missing is a \emph{mutation operator}: a procedure that derives new candidates from the survivors rather than proposing from scratch.
Adding mutation transforms the GES loop into the full generate→evaluate→select→\textbf{mutate} evolutionary loop of Chapter~\ref{chapter:EvolutionaryFrameworks}.

The student who understands the GES loop and its failure modes — why candidates are duplicated, why the evaluator can be fooled, why the proposer collapses to a single good candidate — is ready for the evolutionary framework.
The pivot is not a new topic; it is an addition to a loop already understood.

\section*{Further Reading}

\begin{small}
\begin{enumerate}
\item Hong, S., et al., ``MetaGPT: Meta Programming for a Multi-Agent Collaborative Framework,'' \emph{ICLR}, 2024 — a framework for structured multi-agent communication.
\item Wu, Q., et al., ``AutoGen: Enabling Next-Gen LLM Applications via Multi-Agent Conversation,'' \emph{arXiv}, 2023 — a widely used orchestration library for multi-agent conversations.
\item Chase, H., \emph{LangChain Documentation} — a survey of pipeline and memory patterns for LLM applications; available at \href{https://python.langchain.com}{python.langchain.com}.
\item \emph{A note on the ecosystem:} the multi-agent framework landscape has fragmented rapidly, so treat the AutoGen and LangChain references above as historical anchors for the conversation and pipeline patterns rather than as current installation targets. As a reasonable present-day starting point aligned with the course's default stack, Claude's agent SDK; but the field moves quickly, so check current adoption before committing to any framework (LangGraph, Swarm, and others are also in active use).
\item Workshop materials, \code{resources/claude-cli-workshop/} Lesson~04 — sub-agents, agent teams, and headless operation with Claude Code.
\item Anthropic, ``Building Effective Agents,'' \emph{Anthropic Engineering}, 2024 — the composable workflow-pattern vocabulary (prompt chaining, routing, parallelization, orchestrator--workers, evaluator--optimizer) that names the patterns of this chapter.
\item Anthropic, ``How We Built Our Multi-Agent Research System,'' \emph{Anthropic Engineering}, 2025 — a shipped orchestrator--worker system: parallel sub-agents in isolated contexts and the token-cost tradeoff of multi-agent designs.
\item Anthropic, ``Effective Harnesses for Long-Running Agents,'' \emph{Anthropic Engineering}, 2025 — session-spanning harness design and external-memory checkpointing (progress log, task list, git history) for agents that outlast a single context window.
\end{enumerate}
\end{small}
% EVOLVE-BLOCK-END
